\documentclass[../main]{subfiles}
\begin{document}
%\subfile{00_tapa2}
\clearpage\maketitle
\tableofcontents
\newpage

\begin{table}[]
    \centering
    \footnotesize
    \begin{tabular}{|p{3cm}|p{4cm}|p{8cm}|} \hline
        Espacio curricular & Ejes / Bloques Temáticos & Aprendizajes priorizados por COVID-19\\ \hline\hline
        DIBUJO TÉCNICO & ELEMENTOS, INSTRUMENTOS Y NORMAS UTILIZADOS EN DIBUJO TÉCNICO & \begin{itemize} \item Identificación de los elementos de trazado vertical, horizontal y de medición.
\item Representar objetos aplicando las normas de Dibujo Técnico.
\item Comprensión de las Normas I.R.A.M. para Dibujo
Técnico en sus generalidades\end{itemize} \\ \hline
         & EJERCICIOS GEOMÉTRICOS EN EL PLANO, ESCALAS Y ACOTACIONES & \begin{itemize}
             \item Desarrollar técnicas básicas de dibujo en la dimensión del plano. 
             \item Reconocimiento y diagramación de formas geométricas simples en el plano de representación. 
             \item Interpretación y utilización de escalas gráficas. 
             \item Reconocimiento de elementos de acotación
\end{itemize}
\\ \hline
 & FORMAS DE REPRESENTACIÓN DE CUERPOS Y CROQUIZADO & \begin{itemize} \item Representar cuerpos de tres dimensiones en el plano,
utilizando diferentes técnicas.
\item Interpretación de un cuerpo volumétrico.
\item Comprensión de técnicas de representación en el plano.
\item Reconocimiento de objetos del mundo real y los imaginados
mediante proyecciones en el plano\end{itemize}\\ \hline
    \end{tabular}
\end{table}

\end{document}
